\documentclass[../../main.tex]{subfiles}% 注意这里的文件路径不能用 ./main.tex ,否则用latexmk编译子文件会报错
\graphicspath{{\subfix{./image/}}} % 指定图片目录,后续可以直接使用图片文件名
% 注意这里的文件路径不能用 ../../image/ ,否则用latexmk编译子文件会报错

% 例如:
% \begin{figure}[H]
% \centering
% \includegraphics[scale=0.3]{图.png}
% \caption{}
% \label{figure:图}
% \end{figure}
% 注意:上述\label{}一定要放在\caption{}之后,否则引用图片序号会只会显示??.

\begin{document}

\section{线性算子的概念}

\subsection{线性算子和线性泛函的定义}

\begin{definition}[线性算子]
设$\mathscr{X},\mathscr{Y}$是两个线性空间，$D$是$\mathscr{X}$的一个线性子空间。$T:D\to\mathscr{Y}$是一种映射，$D$称为$T$的\textbf{定义域}，有时记作$D(T)$。将$R(T)=\{Tx\mid \forall x\in D\}$称为$T$的\textbf{值域}。如果
\begin{align*}
T(\alpha x+\beta y)=\alpha Tx+\beta Ty\quad (\forall x,y\in D,\forall \alpha,\beta\in\mathbb{K}),
\end{align*}
那么称$T$是一个\textbf{线性算子}。
\end{definition}
\begin{remark}
若$T$是线性空间$\mathscr{X}$到线性空间$\mathscr{Y}$的线性算子,则$T\theta=\theta$.事实上,若$T(\theta)=y\neq \theta$,则$T(t\theta)=tT(\theta)=y,\forall t\in \mathbb{K}$,故$y=\theta$.
\end{remark}

\begin{proposition}
\begin{enumerate}[(1)]
\item 设$\mathscr{X}=\mathbb{R}^n,\mathscr{Y}=\mathbb{R}^m,T=(t_{ij})_{m\times n}$。如果
\begin{align*}
x\mapsto Tx=\left( \sum\limits_{j=1}^n t_{ij}x_j \right)_{i=1}^m\quad (\forall x=(x_1,x_2,\cdots,x_n)\in\mathbb{R}^n),
\end{align*}
那么$T$是一个线性算子。

\item 设$\mathscr{X}=\mathscr{Y}=C^\infty(\overline{\Omega})$，又设微分多项式
\begin{align*}
P(\partial_x)=\sum\limits_{|\alpha|\leqslant m} a_\alpha(x)\partial_x^\alpha\quad (a_\alpha(x)\in C^\infty(\overline{\Omega})).
\end{align*}
如果$T:u(x)\mapsto P(\partial_x)u(x)(\forall u\in\mathscr{X})$，那么$T$便是一个$\mathscr{X}$到$\mathscr{Y}$的线性算子。

若$\mathscr{X}=\mathscr{Y}=L^2(\Omega),D(T)=C^m(\overline{\Omega})$，则上面定义的算子$T$也是线性的。

\item 设$\mathscr{X}=L^1(-\infty,\infty),\mathscr{Y}=L^\infty(-\infty,\infty)$，若规定
\begin{align*}
T:u(x)\mapsto \int_{-\infty}^\infty e^{\mathrm{i}\xi\cdot x}u(x)\mathrm{d}x\quad (\forall u\in\mathscr{X}),
\end{align*}
那么$T$是一个$\mathscr{X}$到$\mathscr{Y}$的线性算子。
\end{enumerate}
\end{proposition}
\begin{proof}

\end{proof}

\begin{definition}[线性泛函]
取值于实数（复数）的线性算子称为\textbf{实（复）线性泛函}，记作$f(x)$或$\langle f,x\rangle$（即线性函数）。
\end{definition}

\begin{proposition}
\begin{enumerate}[(1)]
\item 设$\mathscr{X}=C(\overline{\Omega})$，若规定
\begin{align*}
f(x)\triangleq \int_\Omega x(\xi)\mathrm{d}\xi\quad (\forall x\in\mathscr{X}),
\end{align*}
则$f$是一个线性泛函，但$x(\xi)\mapsto \int_\Omega x^2(\xi)\mathrm{d}\xi$却不是线性泛函。

\item 设$\mathscr{X}=C^\infty(\Omega)$，若对某个指标$\alpha$及$\xi_0\in\Omega$规定
\begin{align*}
f(u)=\partial^\alpha u(\xi_0)\quad (\forall u\in\mathscr{X}),
\end{align*}
则$f$是$C^\infty(\Omega)$上的一个线性泛函。
\end{enumerate}
\end{proposition}
\begin{proof}

\end{proof}



\subsection{线性算子的连续性和有界性}

\begin{definition}
设$\mathscr{X},\mathscr{Y}$是$B^*$空间，$D(T)\subseteq\mathscr{X}$，称线性算子$T:D(T)\to\mathscr{Y}$在$x_0\in D(T)$是\textbf{连续的}，如果
\begin{align*}
x_n\in D(T),\ x_n\to x_0\implies Tx_n\to Tx_0.
\end{align*}
\end{definition}

\begin{proposition}\label{proposition:线性算子在零元处连续等价于全空间连续}
线性算子$T$在$D(T)$内处处连续当且仅当它在$x=\theta$(零向量)处连续。
\end{proposition}
\begin{proof}
若$T$在$\theta$处连续，那么对$\forall x_n,x_0\in D(T),x_n\to x_0$有
\begin{align*}
x_n-x_0\to\theta\implies Tx_n-Tx_0=T(x_n-x_0)\to T\theta=\theta.
\end{align*}

\end{proof}

\begin{definition}
设$\mathscr{X},\mathscr{Y}$都是$B^*$空间，称线性算子$T:\mathscr{X}\to\mathscr{Y}$是\textbf{有界的}，如果有常数$M\geqslant 0$，使得
\begin{align*}
\|Tx\|_{\mathscr{Y}}\leqslant M\|x\|_{\mathscr{X}}\quad (\forall x\in\mathscr{X}).
\end{align*}
\end{definition}

\begin{proposition}\label{proposition:线性算子有界当且仅当连续}
设$\mathscr{X},\mathscr{Y}$都是$B^*$空间，线性算子$T$连续当且仅当$T$有界。
\end{proposition}
\begin{proof}
充分性显然。下证必要性。若不然，则对$\forall n\in \mathbb{N}$,$\exists x_n\in\mathscr{X}$，使得
\begin{align*}
\|Tx_n\|>n\|x_n\|.
\end{align*}
令$y_n=\frac{x_n}{n\|x_n\|}$，便有$\|Ty_n\|>1$。但$y_n\to\theta(n\to\infty)$，便与$T$的连续性矛盾！

\end{proof}

\begin{definition}
用$\mathscr{L}(\mathscr{X},\mathscr{Y})$表示一切由线性空间$\mathscr{X}$到线性空间$\mathscr{Y}$的有界线性算子的全体，并规定
\begin{align*}
\|T\|=\sup_{x\in\mathscr{X}\setminus\{\theta\}}\frac{\|Tx\|}{\|x\|}=\sup_{\|x\|=1}\|Tx\|
\end{align*}
为$T\in\mathscr{L}(\mathscr{X},\mathscr{Y})$的\textbf{范数},其中$\theta$是零向量。特别用$\mathscr{L}(\mathscr{X})$表示$\mathscr{L}(\mathscr{X},\mathscr{X})$，用$\mathscr{X}^*$表示$\mathscr{L}(\mathscr{X},\mathbb{K})$，即$\mathscr{X}^*$表示$\mathscr{X}$上的有界线性泛函全体,其中$\mathbb{K}=\mathbb{C}\text{或}\mathbb{R}$.
\end{definition}
\begin{remark}
对$\forall x\in \mathscr{X} \setminus \{\theta\}$，令$u=\frac{x}{\|x\|}$，则$\|u\|=1$。从而
\begin{align*}
\frac{\|Tx\|}{\|x\|}=\frac{\|T(\|x\| u)\|}{\|x\|}=\frac{\|x\| \cdot \|Tu\|}{\|x\|}=\|Tu\|.
\end{align*}
于是$\sup\limits_{x\in \mathscr{X} \setminus \{\theta\}}\frac{\|Tx\|}{\|x\|}\leqslant \sup\limits_{\|u\|=1}\|Tu\|.$又显然有
\begin{align*}
\|Tx\|=\frac{\|Tx\|}{\|x\|}\ (\forall x\in \mathscr{X}且\|x\|=1)\Longrightarrow \sup\limits_{\|x\|=1}\|Tx\|\leqslant \sup\limits_{x\in \mathscr{X} \setminus \{\theta\}}\frac{\|Tx\|}{\|x\|}.
\end{align*}
故$\sup_{x\in\mathscr{X}\setminus\{\theta\}}\frac{\|Tx\|}{\|x\|}=\sup_{\|x\|=1}\|Tx\|.$
\end{remark}

\begin{proposition}\label{proposition:有界线性算子的基本不等式}
设$T$是线性空间$\mathscr{X}$到线性空间$\mathscr{Y}$的一个有界线性算子,即$T\in \mathscr{L}(\mathscr{X},\mathscr{Y})$,则
\begin{align*}
\left\| Tx \right\| \leqslant \left\| T \right\| \left\| x \right\| ,\quad \forall x\in \mathscr{X} .
\end{align*}
\end{proposition}
\begin{proof}
由$T$的范数的定义知
\begin{align*}
\left\| T \right\| =\underset{x\in \mathscr{X} \backslash \left\{ \theta \right\}}{\mathrm{sup}}\frac{\left\| Tx \right\|}{\left\| x \right\|}\geqslant \frac{\left\| Tx \right\|}{\left\| x \right\|}\left( \forall x\in \mathscr{X} \backslash \left\{ \theta \right\} \right) \Longrightarrow \left\| Tx \right\| \leqslant \left\| T \right\| \left\| x \right\| \left( \forall x\in \mathscr{X} \backslash \left\{ \theta \right\} \right) .
\end{align*}
当$x=\theta$时,显然有$\left\| Tx \right\| =\left\| T \right\| \left\| x \right\| =0.$故
\begin{align*}
\left\| Tx \right\| \leqslant \left\| T \right\| \left\| x \right\| ,\quad \forall x\in \mathscr{X} .
\end{align*}

\end{proof}

\begin{theorem}\label{theorem:有界线性算子全体构成的空间完备}
设$\mathscr{X}$是$B^*$空间，$\mathscr{Y}$是$B$空间，若在$\mathscr{L}(\mathscr{X},\mathscr{Y})$上规定线性运算：
\begin{align*}
(\alpha_1T_1+\alpha_2T_2)(x)=\alpha_1T_1x+\alpha_2T_2x\quad (\forall x\in\mathscr{X}),
\end{align*}
其中$\alpha_1,\alpha_2\in\mathbb{K},T_1,T_2\in\mathscr{L}(\mathscr{X},\mathscr{Y})$，则$\mathscr{L}(\mathscr{X},\mathscr{Y})$按$\|T\|$构成一个$B$空间。
\end{theorem}
\begin{proof}
显然$\mathscr{L}(\mathscr{X},\mathscr{Y})$是一个线性空间，下证$\|T\|$是范数：
\begin{align*}
\|T\|\geqslant 0,\quad \|T\|=0\iff Tx=\theta(\forall x\in\mathscr{X})\iff T=\theta,
\end{align*}
\begin{align*}
\|T_1+T_2\|=\sup_{\|x\|=1}\|T_1x+T_2x\|
\leqslant \sup_{\|x\|=1}\|T_1x\|+\sup_{\|x\|=1}\|T_2x\|
=\|T_1\|+\|T_2\|,
\end{align*}
\begin{align*}
\|\alpha T\|=\sup_{\|x\|=1}\|\alpha Tx\|=|\alpha|\sup_{\|x\|=1}\|Tx\|=|\alpha|\|T\|.
\end{align*}

再证完备性。设$\{T_n\}_1^\infty$是一个基本列，则$\forall \varepsilon>0,\exists N=N(\varepsilon)$，使得
\begin{align*}
\left\| T_{n+p}-T_n \right\| =\underset{x\in \mathscr{X} \backslash \left\{ \theta \right\}}{\mathrm{sup}}\frac{\left\| T_{n+p}x-T_nx \right\|}{\left\| x \right\|}<\varepsilon \quad (\forall n>N,\forall p\in \mathbb{N} ).
\end{align*}
对$\forall x\in\mathscr{X}$,固定$x$,有
\begin{align*}
\left\| T_{n+p}x-T_nx \right\| \leqslant \varepsilon \left\| x \right\| \quad (\forall n>N,\forall p\in \mathbb{N} ).
\end{align*}
因此$\{T_nx\}_{n=1}^{\infty}$是$\mathscr{Y}$中的基本列,于是由$\mathscr{Y}$是$B$空间知$T_nx\to y\in\mathscr{Y}(n\to\infty)$。记此$y=Tx$，则由$x$的任意性可定义
\begin{align*}
T:\mathscr{X}\to \mathscr{Y},\,\,x\mapsto Tx.
\end{align*}
则
\begin{align*}
\left\| T_n-T \right\| =\underset{\left\| x \right\| =1}{\mathrm{sup}}\left\| T_nx-Tx \right\| \rightarrow 0\quad \left( n\rightarrow \infty \right) .
\end{align*}
下面我们只需证$T\in\mathscr{L}(\mathscr{X},\mathscr{Y})$。不难看出$T$是线性的，再证其有界。事实上，$\exists N\in\mathbb{N}$，使得
\begin{align*}
\|Tx\|=\|y\|\leqslant \|T_Nx\|+1\leqslant \|T_N\|+1=(\|T_N\|+1)\|x\|\quad (\forall x\in\mathscr{X},\|x\|=1).
\end{align*}
即得$\|T\|\leqslant \|T_N\|+1$.因此$\|T\|$是有界的.

\end{proof}

\begin{proposition}
设$T$是有穷维$B^*$空间$\mathscr{X}$到$\mathscr{Y}$的线性映射，则$T$必是连续的。
\end{proposition}
\begin{proof}
分别取定空间$\mathscr{X}$和$\mathscr{Y}$的两组基,则$T$可以通过矩阵$(t_{ij})$表示出来，而同一个有穷维空间的任意两个范数等价。不妨取$\mathscr{X}=\mathbb{K}^n,\mathscr{Y}=\mathbb{K}^m$，利用Cauchy-Schwarz不等式便有
\begin{align*}
\|Tx\|=\left( \sum\limits_{i=1}^m \left| \sum\limits_{j=1}^n t_{ij}x_j \right|^2 \right)^{\frac{1}{2}}
\leqslant \left( \sum\limits_{i=1}^m \sum\limits_{j=1}^n |t_{ij}|^2 \cdot \sum\limits_{j=1}^n |x_j|^2 \right)^{\frac{1}{2}}
=\left( \sum\limits_{i=1}^m \sum\limits_{j=1}^n |t_{ij}|^2 \right)^{\frac{1}{2}} \|x\|.
\end{align*}
即得$\|Tx\|\leqslant \left( \sum\limits_{i=1}^m \sum\limits_{j=1}^n |t_{ij}|^2 \right)^{\frac{1}{2}} \|x\|$.

\end{proof}

\begin{theorem}
Hilbert空间$\mathscr{X}$上的正交投影算子。设$M$是$\mathscr{X}$的一个闭线性子空间，依\hyperref[corollary:泛函分析--正交分解定理]{正交分解定理}，$\forall x\in\mathscr{X}$，存在唯一的分解
\begin{align*}
x=y+z,
\end{align*}
其中$y\in M,z\in M^\perp$。对应$x\mapsto y$称作由$\mathscr{X}$到$M$的\textbf{正交投影算子}，记作$P_M$。在不强调子空间$M$时，我们省略$M$而简记为$P$。而且$P$还是一个连续线性算子，并且如果$M\neq\{\theta\}$，那么$\|P\|=1$。
\end{theorem}
\begin{proof}
先证线性。设$x_i=Px_i+z_i(i=1,2)$，其中$z_i\in M^\perp$。这时，
\begin{align*}
\alpha_1x_1+\alpha_2x_2=(\alpha_1Px_1+\alpha_2Px_2)+(\alpha_1z_1+\alpha_2z_2)\quad (\forall \alpha_1,\alpha_2\in\mathbb{K}).
\end{align*}
因为$\alpha_1z_1+\alpha_2z_2\in M^\perp$，而$\alpha_1Px_1+\alpha_2Px_2\in M$，所以
\begin{align*}
P(\alpha_1x_1+\alpha_2x_2)=\alpha_1Px_1+\alpha_2Px_2.
\end{align*}
即得$P$是线性算子。

其次证连续，由\rrefpro{proposition:正交的基本性质}{proposition:正交的基本性质-2}可知$\|Px\|^2=\|x\|^2-\|z\|^2\leqslant\|x\|^2$。因此，$\|Px\|\leqslant\|x\|$，或者$\|P\|\leqslant1$.

最后，当$M\neq\{\theta\}$时，任取$x\in M\setminus\{\theta\}$，便有$Px=x$,于是$\|Px\|=\|x\|$，从而$\|P\|=1$.

\end{proof}






(本节各题中, $\mathscr{X},\mathscr{Y}$ 均指 $B$空间)

\begin{example}
求证: $T \in \mathscr{L}(\mathscr{X},\mathscr{Y})$ 的充要条件是 $T$ 为线性算子, 并将 $\mathscr{X}$ 中的有界集映为 $\mathscr{Y}$ 中的有界集.
\end{example}

\begin{example}
设 $A \in \mathscr{L}(\mathscr{X},\mathscr{Y})$, 求证:
\begin{enumerate}[(1)]
\item $\|A\| = \sup_{\|x\| \leqslant 1} \|Ax\|$;
\item $\|A\| = \sup_{\|x\| < 1} \|Ax\|$.
\end{enumerate}
\end{example}

\begin{example}
设 $f \in \mathscr{L}(\mathscr{X},\mathbb{R})$, 求证:
\begin{enumerate}[(1)]
\item $\|f\| = \sup_{\|x\| = 1} f(x)$;
\item $\sup_{\|x\| < \delta} f(x) = \delta\|f\| \quad (\forall \delta > 0)$.
\end{enumerate}
\end{example}

\begin{example}
设 $y(t) \in C[0,1]$, 定义 $C[0,1]$ 上的泛函
\begin{align*}
f(x) = \int_{0}^{1} x(t)y(t) \mathrm{d}t \quad (\forall x \in C[0,1]),
\end{align*}
求 $\|f\|$.
\end{example}

\begin{example}
设 $f$ 是 $\mathscr{X}$ 上的非零有界线性泛函, 令
\begin{align*}
d = \inf \{\|x\| \mid f(x) = 1, x \in \mathscr{X}\},
\end{align*}
求证: $\|f\| = \frac{1}{d}$.
\end{example}

\begin{example}
设 $f \in \mathscr{X}^*$, 求证: $\forall \varepsilon > 0, \exists x_0 \in \mathscr{X}$, 使得 $f(x_0) = \|f\|$, 且 $\|x_0\| < 1 + \varepsilon$.
\end{example}

\begin{proposition}\label{proposition:泛函分析--命题1.1231.31.3}
设 $T: \mathscr{X} \to \mathscr{Y}$ 是线性的, 令
\begin{align*}
N(T) \triangleq \{x \in \mathscr{X} \mid Tx = \theta\}.
\end{align*}
\begin{enumerate}[(1)]
\item\label{proposition:泛函分析--命题1.1231.31.3-1} 若 $T \in \mathscr{L}(\mathscr{X},\mathscr{Y})$, 求证: $N(T)$ 是 $\mathscr{X}$ 的闭线性子空间.

\item\label{proposition:泛函分析--命题1.1231.31.3-2} 问 $N(T)$ 是 $\mathscr{X}$ 的闭线性子空间能否推出 $T \in \mathscr{L}(\mathscr{X},\mathscr{Y})$?

\item\label{proposition:泛函分析--命题1.1231.31.3-3} 若 $f$ 是线性泛函, 求证:
\begin{align*}
f \in \mathscr{X}^* \iff N(f) \text{ 是闭线性子空间}.
\end{align*}
\end{enumerate}
\end{proposition}

\begin{example}
设 $f$ 是 $\mathscr{X}$ 上的线性泛函, 记
\begin{align*}
H_f^\lambda \triangleq \{x \in \mathscr{X} \mid f(x) = \lambda\} \quad (\forall \lambda \in \mathbb{K}).
\end{align*}
如果 $f \in \mathscr{X}^*$, 并且 $\|f\| = 1$, 求证:
\begin{enumerate}[(1)]
\item $|f(x)| = \inf \{\|x - z\| \mid z \in H_f^0\} \quad (\forall x \in \mathscr{X})$;
\item $\forall \lambda \in \mathbb{K}, H_f^\lambda$ 上的任一点 $x$ 到 $H_f^0$ 的距离都等于 $|\lambda|$.
\end{enumerate}
并对 $\mathscr{X} = \mathbb{R}^2, \mathbb{K} = \mathbb{R}$ 情形解释 (1) 和 (2) 的几何意义.
\end{example}

\begin{proposition}\label{proposition:泛函分析--例题2.1.9}
设 $\mathscr{X}$ 是实 $B^*$ 空间, $f$ 是 $\mathscr{X}$ 上的非零实值线性泛函, 求证: 不存在开球 $B(x_0, \delta)$, 使得 $f(x_0)$ 是 $f(x)$ 在 $B(x_0, \delta)$ 中的极大值或极小值.
\end{proposition}








\end{document}